\cvsection{Open-Source Projects}
\leftright{\textbf{LLM-for-Zotero}: AI agent system for personal academic library management (\link{https://github.com/yilewang/llm-for-zotero}{GitHub})}{\textit{TypeScript}}
\begin{itemize}
  \item \textbf{Architected an end-to-end agent harness} that turns Zotero from a reference manager into a research environment where AI agents work over users' scientific libraries, call tools, and autonomously run multi-step research workflows.
  \item \textbf{Designed the tool-use and orchestration layer}, exposing structured operations for literature retrieval, PDF and figure analysis, multi-paper comparison, scholarly search, metadata and collection management, and local file operations.
  \item \textbf{Engineered long-horizon context and workflow systems}: source-coverage tracking, context compaction, reusable agent skills, caching, a database layer, and human-in-the-loop confirmation and undo for autonomous actions.
  \item \textbf{Built a model-agnostic agent runtime} supporting OpenAI-compatible APIs, local models, Codex App Server, Claude Code, and other LLM providers behind one unified experience inside Zotero.
  \item \textbf{Grew the project to 3,000+ GitHub stars and 600K+ downloads}, iterating on a widely adopted AI product with a fast-growing research community.
\end{itemize}
\entrygap
\leftright{\textbf{Human Brain-like Vision Model} (\link{https://github.com/ZitongLu1996/ReAlnet}{GitHub})}{\textit{PyTorch}}
\begin{itemize}
  \item \textbf{Developed a brain-guided deep learning framework} that incorporates human neural representations into DCNN training, improving alignment between model representations and human brain activity.
  \item \textbf{Designed a personalized, layer-wise alignment pipeline} using Representational Similarity Analysis (RSA), adapting models to individual neural and behavioral response patterns.
  \item Implemented multi-layer feature extraction, neural representation matching, and model optimization in \textbf{PyTorch} for large-scale computational neuroscience experiments.
\end{itemize}
\entrygap
\leftright{\textbf{Dynamical Network Analysis \& Statistical Toolbox} (\link{https://github.com/yilewang/dynamics_net_sys_class/tree/main}{GitHub})}{\textit{Python}}
\begin{itemize}
  \item Built a Python toolbox for modeling and analyzing complex dynamical networks, \textbf{integrating dynamical systems theory, network science, and graph theory}.
  \item \textbf{Implemented and visualized canonical models}, including SIR epidemic dynamics, network diffusion, and Erdős--Rényi random graphs, with reusable simulation and analysis workflows.
  \item Created interactive notebooks that connect mathematical models to numerical simulation, network structure, and emergent system dynamics.
\end{itemize}
